% CP-VI-0034
% target_chapter: VI/26 — Finite-Type and Noetherian Morphisms
% source_unit_id: IMP-DL-D27460A55B-F2R000465-7C515FA1
% source: Downloads/theory-of-algebraic-geometry-5.tex L465-L483
% solution_provenance: CANONICAL_AUTHORED_SOLUTION

\textbf{Problem.}
We say a scheme \(X\) is \emph{locally Noetherian} if it can be covered by affine open subsets
\[
\operatorname{Spec} A_i
\]
with \(A_i\) a Noetherian ring.

We say a scheme \(X\) is \emph{Noetherian} if it can be covered by a finite number of affine open subsets
\[
\operatorname{Spec} A_i
\]
with \(A_i\) Noetherian.

Show that a scheme \(X\) is locally Noetherian if and only if for every open affine subset
\[
U=\operatorname{Spec} A,
\]
the ring \(A\) is Noetherian.

\subsection*{Canonical solution}
Suppose first that $X$ is locally Noetherian and let
$U=\operatorname{Spec} A\subseteq X$ be any affine open.
For each $x\in U$, choose an affine Noetherian neighbourhood
$V=\operatorname{Spec} C$ of $x$ in $X$. A standard
simultaneous-distinguished-open refinement of $U\cap V$
(as proved in CP-VI-0032) gives elements $f\in A$ and $g\in C$
such that
\[
x\in D_U(f)=D_V(g)\subseteq U\cap V.
\]
Since $C$ is Noetherian, the localization $C_g$ is Noetherian;
hence $A_f\cong C_g$ is Noetherian. These principal opens $D_U(f)$
cover $U$. Quasi-compactness gives finitely many
$D_U(f_1),\ldots,D_U(f_r)$ covering $U$, so
$(f_1,\ldots,f_r)=A$ and each $A_{f_j}$ is Noetherian.
The distinguished-cover criterion for Noetherian rings (proved by
clearing denominators in the finitely many local generators of any
ideal) now implies that $A$ is Noetherian.

Conversely, if every affine open $\operatorname{Spec} A\subseteq X$ has $A$ Noetherian, then any affine open cover of $X$ is a cover by spectra of Noetherian rings. Hence $X$ is locally Noetherian.
